\documentclass[10pt,a4paper]{amsart}
\usepackage{amsmath,amssymb,amsthm,graphicx}
% Retained minimal preamble for the Stacks Project chapter varieties.tex.
% Verbatim excerpt of the chapter preamble: Xy-pic setup (native preamble lines 14-18), the native theorem
% declarations with the mathematical macros (native preamble lines 154-217), and the project own no-op
% definitions of the reference/slogan/history display environments (native wrapper.tex lines 1-3) that the
% counted statement uses for its slogan. Copied mechanically; nothing added, removed or redefined.
% For commutative diagrams we use Xy-pic
\usepackage[all]{xy}

% We use 2cell for 2-commutative diagrams.
\xyoption{2cell}
% Theorem environments.
%
\theoremstyle{plain}
\newtheorem{theorem}[subsection]{Theorem}
\newtheorem{proposition}[subsection]{Proposition}
\newtheorem{lemma}[subsection]{Lemma}

\theoremstyle{definition}
\newtheorem{definition}[subsection]{Definition}
\newtheorem{example}[subsection]{Example}
\newtheorem{exercise}[subsection]{Exercise}
\newtheorem{situation}[subsection]{Situation}

\theoremstyle{remark}
\newtheorem{remark}[subsection]{Remark}
\newtheorem{remarks}[subsection]{Remarks}

\numberwithin{equation}{subsection}

% Macros
%
\def\lim{\mathop{\mathrm{lim}}\nolimits}
\def\colim{\mathop{\mathrm{colim}}\nolimits}
\def\Spec{\mathop{\mathrm{Spec}}}
\def\Hom{\mathop{\mathrm{Hom}}\nolimits}
\def\Ext{\mathop{\mathrm{Ext}}\nolimits}
\def\SheafHom{\mathop{\mathcal{H}\!\mathit{om}}\nolimits}
\def\SheafExt{\mathop{\mathcal{E}\!\mathit{xt}}\nolimits}
\def\Sch{\mathit{Sch}}
\def\Mor{\mathop{\mathrm{Mor}}\nolimits}
\def\Ob{\mathop{\mathrm{Ob}}\nolimits}
\def\Sh{\mathop{\mathit{Sh}}\nolimits}
\def\NL{\mathop{N\!L}\nolimits}
\def\CH{\mathop{\mathrm{CH}}\nolimits}
\def\proetale{{pro\text{-}\acute{e}tale}}
\def\etale{{\acute{e}tale}}
\def\QCoh{\mathit{QCoh}}
\def\Ker{\mathop{\mathrm{Ker}}}
\def\Im{\mathop{\mathrm{Im}}}
\def\Coker{\mathop{\mathrm{Coker}}}
\def\Coim{\mathop{\mathrm{Coim}}}

% Boxtimes
%
\DeclareMathSymbol{\boxtimes}{\mathbin}{AMSa}{"02}

%
% Macros for moduli stacks/spaces
%
\def\QCohstack{\mathcal{QC}\!\mathit{oh}}
\def\Cohstack{\mathcal{C}\!\mathit{oh}}
\def\Spacesstack{\mathcal{S}\!\mathit{paces}}
\def\Quotfunctor{\mathrm{Quot}}
\def\Hilbfunctor{\mathrm{Hilb}}
\def\Curvesstack{\mathcal{C}\!\mathit{urves}}
\def\Polarizedstack{\mathcal{P}\!\mathit{olarized}}
\def\Complexesstack{\mathcal{C}\!\mathit{omplexes}}
% \Pic is the operator that assigns to X its picard group, usage \Pic(X)
% \Picardstack_{X/B} denotes the Picard stack of X over B
% \Picardfunctor_{X/B} denotes the Picard functor of X over B
\def\Pic{\mathop{\mathrm{Pic}}\nolimits}
\def\Picardstack{\mathcal{P}\!\mathit{ic}}
\def\Picardfunctor{\mathrm{Pic}}
\def\Deformationcategory{\mathcal{D}\!\mathit{ef}}
\newenvironment{reference}{}{}
\newenvironment{slogan}{}{}
\newenvironment{history}{}{}
\usepackage{cleveref}
\crefname{originallemma}{original lemma}{original lemmas}
\Crefname{originallemma}{Original lemma}{Original lemmas}
\makeatletter
\newlabel{lemma-create-globally-generated}{{09N5}{source}{}{sourcepointer-1}{}}
\newlabel{lemma-create-globally-generated@cref}{{[originallemma][0][]09N5}{[1][0][]source}{}{}{}}
\makeatother
\title{Existence of a globally generated invertible sheaf with prescribed nonvanishing locus on an integral separated scheme}
\author{The Stacks Project Authors}
\date{stacks-project:varieties.tex:local-snapshot-20261011; GNU FDL 1.2 or later}
\begin{document}
\maketitle
\noindent\textit{Editorial note.} All delivered mathematical text is the author's own, copied verbatim from the pinned Stacks Project chapter file varieties.tex; no statement or proof step was rewritten, reordered or removed.\par\medskip
\noindent\textit{Editorial note.} The counted span cites one cross-chapter label, recorded in assembly.external\_references with its canonical tag as a navigation pointer.\par\medskip
\noindent\textit{Editorial note.} The proof cites no bibliography item inside the counted span.\par\medskip
\noindent\textit{Allowed standard prerequisites (editorial).}\begin{itemize}
\item Existence of an element whose inverse lies in the Jacobson radicals of the local rings at the missing points. Stacks Project, Varieties, Lemma lemma-create-globally-generated (tag 09N5). Supplies the element f from which the cocycle for the line bundle is built.
\end{itemize}\medskip
\ifdefined\hypertarget\hypertarget{sourcepointer-1}{}\fi
\noindent\textit{Source reference.} Original lemma 09N5: sources/stacks-chow/tags.txt, source lines 12203--12203. Cross-chapter standard prerequisite cited by the author inside the counted span; it is used as a pointer and its own proof is not part of this item.\par\medskip
\begin{lemma}
\label{lemma-find-globally-generated}
Let $X$ be an integral separated scheme. Let $U \subset X$ be a nonempty
affine open such that $X \setminus U$ is a finite set of points
$x_1, \ldots, x_r$ with $\mathcal{O}_{X, x_i}$ Noetherian of dimension $1$.
Then there exists a globally generated invertible $\mathcal{O}_X$-module
$\mathcal{L}$ and a section $s$ such that $U = X_s$.
\end{lemma}

\begin{proof}
Say $U = \Spec(A)$ and let $K$ be the function field of $X$.
Write $B_i = \mathcal{O}_{X, x_i}$ and $\mathfrak m_i = \mathfrak m_{x_i}$.
Since $x_i \not \in U$ we see that the open
$U \times_X \Spec(B_i)$ of $\Spec(B_i)$ has only one point, i.e.,
$U \times_X \Spec(B_i) = \Spec(K)$.
Since $X$ is separated, we find that $\Spec(K)$ is a closed subscheme
of $U \times \Spec(B_i)$, i.e., the map $A \otimes B_i \to K$ is a surjection.
By Lemma \ref{lemma-create-globally-generated} we can find a nonzero
$f \in A$ such that $f^{-1} \in \mathfrak m_i$ for $i = 1, \ldots, r$.
Pick opens $x_i \in U_i \subset X$ such that $f^{-1} \in \mathcal{O}(U_i)$.
Then
$$
\mathcal{U} : X = U \cup \bigcup U_i
$$
is an open covering of $X$. Consider the $2$-cocycle with values
in $\mathcal{O}_X^*$ given by $f$ on $U \cap U_i$ and by
$1$ on $U_i \cap U_j$. This defines a line bundle $\mathcal{L}$
with two sections:
\begin{enumerate}
\item a section $s$ defined by $1$ on $U$
and $f^{-1}$ on $U_i$ is as in the statement of the lemma, and
\item a section $t$ defined by $f$ on $U$ and
$1$ on $U_i$.
\end{enumerate}
Note that $X_t \supset U_1 \cup \ldots \cup U_r$.
Hence $s, t$ generate $\mathcal{L}$ and the lemma is proved.
\end{proof}

\end{document}
